\documentclass[11pt, a4paper]{article}

\usepackage[T1]{fontenc}
\usepackage[utf8]{inputenc}
\usepackage{tgpagella}
\usepackage{microtype}
\usepackage[spanish, es-noshorthands]{babel}
\usepackage[table, xcdraw]{xcolor}
\usepackage{booktabs}
\usepackage{tabularx}
\usepackage[margin=2.5cm]{geometry}
\usepackage{fancyhdr}

\pagestyle{fancy}
\fancyhf{}
\setlength{\headheight}{16pt}
\lhead{\textbf{UCB Sede Tarija} | Ing. Mecatrónica}
\chead{Hoja de Ruta de Evaluaciones}
\rhead{Hacia EC2 e Hito 2}

\definecolor{ucbblue}{RGB}{0, 51, 102}

\begin{document}

\begin{center}
    \Large\textbf{\textcolor{ucbblue}{Roadmap de Evaluaciones (Provisional)}}\\[0.2cm]
\end{center}

\textit{Nota: Las fechas presentadas aquí son de carácter provisional y están sujetas a la verificación (Readiness Gate) de los temas de Singularidad y Dinámica[cite: 13].}

\vspace{0.4cm}
\renewcommand{\arraystretch}{1.5}
\begin{tabularx}{\textwidth}{|l|l|X|l|}
\hline
\rowcolor{ucbblue!20} \textbf{Fecha} & \textbf{Actividad} & \textbf{Propósito Pedagógico} & \textbf{Estado} \\ \hline
Mar 29 Sep & Singularidades + LSPB & Cerrar brechas conceptuales del Jacobiano y abrir perfil temporal. & \textit{Programado} \\ \hline
Jue 1 Oct & Dinámica e Integración H2 & Bloque final de instrucción masiva (Ecuación Torque). & \textit{Provisional} \\ \hline
\textbf{Mar 6 Oct} & \textbf{Evaluación EC2 (Individual)} & ¿Puede el estudiante resolver la matemática de base solo? & \textit{Provisional} \\ \hline
\textbf{Jue 8 Oct} & \textbf{Defensa Hito 2 (Grupal)} & ¿Puede el grupo integrar todo en un pipeline de ingeniería coherente? & \textit{Provisional} \\ \hline
\end{tabularx}

\vspace{0.8cm}
\textbf{Naturaleza de la Evaluación[cite: 13]:}
\begin{itemize}
    \item \textbf{EC2:} Enfoque analítico cerrado en lápiz y papel.
    \item \textbf{Hito 2:} Enfoque algorítmico, resolución computacional y defensa de la integración cinemática y dinámica en Python (Notebook).
\end{itemize}

\end{document}
